\documentclass[a4paper,12pt]{article}
\usepackage[10pt]{type1ec}
\usepackage[T2A]{fontenc}
\pdfmapfile{+cm-super-extended-t2a.map}
\usepackage[utf8]{inputenc}
\usepackage[russian]{babel}
\usepackage{amsmath,amssymb,mathtools}
\usepackage{geometry}
\geometry{left=23mm,right=23mm,top=23mm,bottom=24mm}
\usepackage{microtype}
\usepackage{icomma}
\usepackage{tikz}
\usetikzlibrary{arrows.meta,positioning,shapes.geometric,calc}
\usepackage{booktabs,tabularx,array}
\usepackage{enumitem}
\usepackage{fancyhdr}
\usepackage{parskip}
\usepackage{xcolor}
\usepackage{hyperref}
\definecolor{accent}{HTML}{007F81}
\definecolor{darkaccent}{HTML}{00545A}
\definecolor{ink}{HTML}{252A2D}
\definecolor{muted}{HTML}{5C6669}
\color{ink}
\hypersetup{colorlinks=true,linkcolor=darkaccent,urlcolor=darkaccent}
\setlength{\parindent}{0pt}
\setlength{\parskip}{6pt plus 1pt minus 1pt}
\linespread{1.16}
\setlist[itemize]{leftmargin=1.45em,itemsep=2pt,topsep=4pt}
\setlist[enumerate]{leftmargin=1.6em,itemsep=6pt,topsep=4pt}
\newcommand{\term}[1]{\textcolor{darkaccent}{\textbf{#1}}}
\newcommand{\key}[1]{\textcolor{accent}{#1}}
\newcommand{\ruleline}{\par\vspace{3pt}{\color{accent!65}\hrule height 0.45pt}\vspace{8pt}}
\pagestyle{fancy}
\fancyhf{}
\fancyhead[L]{\footnotesize\color{muted} Тригонометрические преобразования}
\fancyhead[R]{\footnotesize\color{muted} Марк Гукин}
\fancyfoot[C]{\small\color{muted}\thepage}
\renewcommand{\headrulewidth}{0.2pt}
\renewcommand{\footrulewidth}{0pt}
\setlength{\headheight}{15pt}
\newcommand{\topic}[1]{\vspace{9pt}{\large\textcolor{darkaccent}{\textbf{#1}}}\par\vspace{2pt}}
\newcommand{\ex}[1]{\par\medskip\textbf{Пример.} #1\par}
\newcommand{\task}[2]{\item \textbf{#1} #2}
\begin{document}
\thispagestyle{empty}
\begin{center}
\vspace*{25mm}
{\Large\color{darkaccent} МАТЕМАТИКА}\par
\vspace{22mm}
{\fontsize{29}{34}\selectfont Чек-лист}\par
\vspace{18mm}
{\fontsize{24}{29}\selectfont\color{accent} Тригонометрические\\[3pt] преобразования}\par
\vspace{19mm}
{\large Подготовка к ЕГЭ по профильной математике}\par
\vfill
{\large Марк Гукин}\par
\vspace{12mm}
{\large 28 октября 2026 г.}\par
\vspace{18mm}
{\normalsize Лёвин Артём Александрович\\[2pt] эксклюзивно для Марка Гукина}\par
\vspace{18mm}
\begin{tikzpicture}[x=1cm,y=1cm]
 \draw[accent,line width=0.9pt] (0,0) .. controls (1.6,0.75) and (2.7,-0.55) .. (4.5,0.08)
  .. controls (6.1,0.65) and (7.25,-0.55) .. (9.2,0.03);
 \fill[darkaccent] (4.5,0.08) circle (1.35pt);
\end{tikzpicture}
\end{center}
\clearpage

\section*{1. Главная идея занятия}
\ruleline
\term{Принцип: приводите тригонометрические выражения к одинаковым аргументам.}
Если в задаче встречаются $\sin 6\alpha$ и $\cos 3\alpha$, стоит заметить, что $6\alpha=2\cdot3\alpha$. Формула двойного угла позволяет получить синус и косинус одного аргумента $3\alpha$.

\topic{Формулы суммы и разности}
Для любых действительных $a,b$:
\begin{align*}
\sin(a+b)&=\sin a\cos b+\cos a\sin b,\\
\sin(a-b)&=\sin a\cos b-\cos a\sin b,\\
\cos(a+b)&=\cos a\cos b-\sin a\sin b,\\
\cos(a-b)&=\cos a\cos b+\sin a\sin b.
\end{align*}
В формулах синуса знак сохраняется; в формулах косинуса знак между произведениями противоположен знаку между углами.

\ex{Найдите $\sin\bigl(\pi+\frac\pi4\bigr)$.}
\begin{align*}
\sin\left(\pi+\frac\pi4\right)
&=\sin\pi\cos\frac\pi4+\cos\pi\sin\frac\pi4\\
&=0\cdot\frac{\sqrt2}{2}+(-1)\cdot\frac{\sqrt2}{2}
=-\frac{\sqrt2}{2}.
\end{align*}
\ex{Найдите $\cos\bigl(2\pi-\frac\pi3\bigr)$.}
\[
\cos\left(2\pi-\frac\pi3\right)
=\cos2\pi\cos\frac\pi3+\sin2\pi\sin\frac\pi3
=\frac12.
\]

\topic{Как восстановить табличные значения}
Для $0^\circ$, $30^\circ$, $45^\circ$, $60^\circ$, $90^\circ$ значения синуса равны $\sqrt n/2$ при $n=0,1,2,3,4$ соответственно. Для косинуса порядок $n$ обратный: $4,3,2,1,0$.
\begin{center}
\small
\renewcommand{\arraystretch}{1.2}
\begin{tabular}{@{}c*{5}{>{\centering\arraybackslash}p{22mm}}@{}}
\toprule
Угол & $0^\circ$ & $30^\circ$ & $45^\circ$ & $60^\circ$ & $90^\circ$\\
\midrule
$\sin$ & $0$ & $\frac12$ & $\frac{\sqrt2}{2}$ & $\frac{\sqrt3}{2}$ & $1$\\
$\cos$ & $1$ & $\frac{\sqrt3}{2}$ & $\frac{\sqrt2}{2}$ & $\frac12$ & $0$\\
\bottomrule
\end{tabular}
\end{center}

\clearpage
\section*{2. Углы, которых нет в таблице}
\ruleline
\term{Первый ход:} разделить числитель угла на кратное знаменателю число и небольшой остаток. Затем упростить дробь и применить формулу суммы или разности.
\begin{align*}
\frac{5\pi}{4}&=\frac{4\pi+\pi}{4}=\pi+\frac\pi4,\\
\frac{7\pi}{3}&=\frac{6\pi+\pi}{3}=2\pi+\frac\pi3,\\
\frac{8\pi}{3}&=\frac{9\pi-\pi}{3}=3\pi-\frac\pi3.
\end{align*}
Здесь используется обычное почленное деление: $\frac{x+y}{n}=\frac xn+\frac yn$ при $n\ne0$.

\ex{Вычислите $\cos\frac{19\pi}{4}$.}
\begin{align*}
\frac{19\pi}{4}&=\frac{16\pi+3\pi}{4}=4\pi+\frac{3\pi}{4},\\
\cos\frac{19\pi}{4}&=\cos\left(4\pi+\frac{3\pi}{4}\right)
=\cos\frac{3\pi}{4}=-\frac{\sqrt2}{2}.
\end{align*}
Для последнего перехода используем период косинуса $2\pi$ и знак косинуса во II четверти. 

\topic{Градусы и радианы}
Один полный оборот равен $360^\circ=2\pi$ радиан, а $180^\circ=\pi$ радиан. Поэтому
\[
\boxed{\quad \varphi_{\text{рад}}=\frac{\varphi_{\text{град}}\,\pi}{180}\quad}
\]
При подстановке числа градусов знак $^\circ$ используется только в обозначении единицы измерения:
\[
270^\circ=\frac{270\pi}{180}=\frac{3\pi}{2},\qquad
\sin270^\circ=\sin\frac{3\pi}{2}=-1.
\]

\topic{Знаки в четвертях}
\begin{center}
\begin{tikzpicture}[scale=0.85]
 \draw[->,thin,muted] (-3.05,0)--(3.1,0) node[right] {$x$};
 \draw[->,thin,muted] (0,-2.0)--(0,2.15) node[above] {$y$};
 \draw[accent,line width=0.75pt] (0,0) circle(1.62);
 \node[align=center] at (1.7,1.15) {\small I\\[-2pt]\scriptsize $\sin>0,\;\cos>0$};
 \node[align=center] at (-1.72,1.15) {\small II\\[-2pt]\scriptsize $\sin>0,\;\cos<0$};
 \node[align=center] at (-1.72,-1.2) {\small III\\[-2pt]\scriptsize $\sin<0,\;\cos<0$};
 \node[align=center] at (1.72,-1.2) {\small IV\\[-2pt]\scriptsize $\sin<0,\;\cos>0$};
\end{tikzpicture}
\end{center}
Квадрат тригонометрической функции может быть известен, но при извлечении корня появятся два знака. Указанный промежуток угла позволяет выбрать нужный.

\clearpage
\section*{3. Формулы двойного угла}
\ruleline
\topic{Синус двойного угла: одна формула}
\[
\boxed{\ \sin 2x=2\sin x\cos x\ }
\qquad\Longleftrightarrow\qquad
\sin x=2\sin\frac x2\cos\frac x2.
\]
\ex{Разложите $\sin 10\alpha$ и $\sin 11\alpha$.}
\[
\sin10\alpha=2\sin5\alpha\cos5\alpha,\qquad
\sin11\alpha=2\sin\frac{11\alpha}{2}\cos\frac{11\alpha}{2}.
\]
Коэффициент $2$ перед произведением остаётся всегда. Аргумент каждого множителя в два раза меньше исходного.

\topic{Косинус двойного угла: три равноправные записи}
\begin{align*}
\cos2x&=\cos^2x-\sin^2x,\\
\cos2x&=2\cos^2x-1,\\
\cos2x&=1-2\sin^2x.
\end{align*}
Если дан синус, удобна последняя запись. Если известен косинус, удобно пользоваться второй. Первая пригодится, когда в выражении присутствуют обе функции.

\topic{Основное тригонометрическое тождество}
\[
\boxed{\sin^2x+\cos^2x=1}
\]
Из него следует:
\[
\cos x=\pm\sqrt{1-\sin^2x},\qquad
\sin x=\pm\sqrt{1-\cos^2x}.
\]
Знак определяется четвертью. Если $\cos^2x=\frac19$, то $\cos x=\pm\frac13$, и до определения четверти выбирать знак нельзя.

\topic{Тангенс}
\[
\operatorname{tg}x=\frac{\sin x}{\cos x},\qquad \cos x\ne0.
\]
Чтобы найти тангенс по известному синусу, сначала восстановите косинус из основного тождества, определите его знак, затем разделите.

\clearpage
\section*{4. Алгоритм: как начать тригонометрическую задачу}
\vspace{8mm}
\begin{center}
\begin{tikzpicture}[
 >=Latex,
 node distance=8.3mm,
 action/.style={draw=accent,rounded corners=3pt,line width=.7pt,align=center,text width=112mm,minimum height=14mm,inner sep=3mm,font=\normalsize},
 terminal/.style={draw=darkaccent,ellipse,line width=.7pt,align=center,text width=78mm,minimum height=12mm,font=\normalsize},
 arr/.style={->,darkaccent,line width=.8pt,shorten >=2pt,shorten <=2pt}
]
\node[terminal] (a) {Перед Вами тригонометрическое выражение};
\node[action,below=of a] (b) {\textbf{Сравните аргументы}\\$2x$ и $x$? Сумма или разность углов?};
\node[action,below=of b] (c) {\textbf{Выберите нужную формулу}\\Двойного угла, суммы или разности};
\node[action,below=of c] (d) {\textbf{Приведите выражения к общему аргументу}\\Упростите дроби, раскройте формулы, сократите только допустимые множители};
\node[action,below=of d] (e) {\textbf{Найдите недостающую функцию, если требуется}\\$\sin^2x+\cos^2x=1$};
\node[action,below=of e] (f) {\textbf{Выберите правильный знак}\\Используйте четверть, заданную промежутком угла};
\node[action,below=of f] (g) {\textbf{Подставьте числа и проверьте вычисления}\\Квадраты, скобки, знаменатели, знаки};
\node[terminal,below=of g] (h) {Получен и проверен ответ};
\foreach \p/\q in {a/b,b/c,c/d,d/e,e/f,f/g,g/h} \draw[arr] (\p)--(\q);
\end{tikzpicture}
\end{center}
\clearpage

\section*{5. Разобранные приёмы}
\ruleline
\topic{Приём 1. Привести аргументы к одинаковому виду}
\ex{Найдите $\dfrac{10\sin6\alpha}{3\cos3\alpha}$, если $\sin3\alpha=0{,}6$ и $\cos3\alpha\ne0$.}
\begin{align*}
\frac{10\sin6\alpha}{3\cos3\alpha}
&=\frac{10\cdot2\sin3\alpha\cos3\alpha}{3\cos3\alpha}\\
&=\frac{20\sin3\alpha}{3}
=\frac{20\cdot0{,}6}{3}=4.
\end{align*}
Сокращение $\cos3\alpha$ допустимо из-за ограничения $\cos3\alpha\ne0$.

\topic{Приём 2. Выбрать подходящую форму косинуса}
\ex{Найдите $24\cos2\alpha$, если $\sin\alpha=-0{,}2$.}
\begin{align*}
\cos2\alpha&=1-2\sin^2\alpha
=1-2(-0{,}2)^2=1-2\cdot0{,}04=0{,}92,\\
24\cos2\alpha&=24\cdot0{,}92=22{,}08.
\end{align*}
\ex{Найдите $2\cos2\alpha$, если $\sin\alpha=-0{,}7$.}
\[
2\cos2\alpha=2\bigl(1-2(-0{,}7)^2\bigr)
=2(1-0{,}98)=0{,}04.
\]
Отрицательный синус возводится в квадрат вместе со знаком, поэтому квадрат положителен.

\topic{Приём 3. Тождество и выбор знака}
\ex{Найдите $3\cos\alpha$, если $\sin\alpha=-\dfrac{2\sqrt2}{3}$, $\dfrac{3\pi}{2}<\alpha<2\pi$.}
\begin{align*}
\cos^2\alpha&=1-\left(-\frac{2\sqrt2}{3}\right)^2
=1-\frac89=\frac19,\\
\cos\alpha&=\pm\frac13.
\end{align*}
В IV четверти косинус положителен: $\cos\alpha=\frac13$. Следовательно, $3\cos\alpha=1$.

\clearpage
\section*{6. Составные задачи: несколько формул подряд}
\ruleline
\topic{Приём 4. Сначала убрать сложный аргумент}
\ex{Найдите $\sin\left(\dfrac{7\pi}{2}-\alpha\right)$, если $\sin\alpha=0{,}8$ и $\dfrac\pi2<\alpha<\pi$.}
\begin{align*}
\sin\left(\frac{7\pi}{2}-\alpha\right)
&=\sin\frac{7\pi}{2}\cos\alpha-\cos\frac{7\pi}{2}\sin\alpha\\
&=(-1)\cos\alpha-0\cdot\sin\alpha=-\cos\alpha.
\end{align*}
Далее найдём косинус по известному синусу:
\[
\cos^2\alpha=1-(0{,}8)^2=0{,}36,\qquad \cos\alpha=\pm0{,}6.
\]
Во II четверти $\cos\alpha<0$, значит $\cos\alpha=-0{,}6$. Поэтому
\[
\boxed{\sin\left(\frac{7\pi}{2}-\alpha\right)=-(-0{,}6)=0{,}6.}
\]

\topic{Приём 5. Из косинуса найти синус двойного угла}
\ex{Известно $\cos\alpha=-0{,}6$ и $\pi<\alpha<\dfrac{3\pi}{2}$. Найдите $\sin2\alpha$.}
\begin{align*}
\sin^2\alpha&=1-(-0{,}6)^2=0{,}64,\\
\sin\alpha&=-0{,}8\quad\text{(III четверть)},\\
\sin2\alpha&=2\sin\alpha\cos\alpha
=2\cdot(-0{,}8)\cdot(-0{,}6)=0{,}96.
\end{align*}

\topic{Типичные ошибки, которые нужно исключить}
\begin{itemize}
\item $\sin(a+b)$ нельзя заменять суммой $\sin a+\sin b$.
\item $\sin2x$ нельзя заменять произведением $2\sin x$.
\item Из $\cos^2x=\frac19$ следуют \emph{оба} варианта $\cos x=\pm\frac13$ до учёта четверти.
\item При подстановке отрицательного числа сохраняйте скобки: $(-0{,}7)^2=0{,}49$.
\item У дробей сохраняйте знаменатель: $1-\frac{144}{169}=\frac{25}{169}$.
\item Сокращать тригонометрический множитель можно только тогда, когда он не равен нулю.
\end{itemize}

\clearpage
\section*{7. Тренировочный блок — путь А}
\ruleline
\smallskip
\textcolor{muted}{10 заданий. Сначала запишите нужную формулу, затем последовательно выполните преобразования.}
\begin{enumerate}[label=\textbf{\arabic*.},leftmargin=2.0em,itemsep=9pt]
\item Вычислите $\displaystyle\sin\left(\pi+\frac\pi6\right)$.
\item Переведите $135^\circ$ в радианы и вычислите $\cos135^\circ$.
\item Вычислите $\displaystyle\cos\left(2\pi-\frac\pi3\right)$.
\item Вычислите $\displaystyle\cos\frac{19\pi}{4}$.
\item Представьте $\sin14\alpha$ в виде произведения по формуле двойного угла.
\item Выразите $\cos8\alpha$ только через $\sin4\alpha$.
\item Найдите $\displaystyle\frac{15\sin8\alpha}{5\cos4\alpha}$, если $\sin4\alpha=0{,}4$ и $\cos4\alpha\ne0$.
\item Найдите $18\cos2\alpha$, если $\sin\alpha=-\dfrac13$.
\item Найдите $\displaystyle\sin\left(\frac{7\pi}{2}-\alpha\right)$, если $\sin\alpha=\dfrac45$ и $\dfrac\pi2<\alpha<\pi$.
\item Найдите $\operatorname{tg}\alpha$, если $\sin\alpha=-\dfrac5{\sqrt{26}}$ и $\pi<\alpha<\dfrac{3\pi}{2}$.
\end{enumerate}
\vspace{8pt}
\topic{Самопроверка перед сдачей}
\begin{itemize}
\item Аргументы тригонометрических функций согласованы с выбранной формулой.
\item Скобки, квадраты и дробные знаменатели сохранены.
\item При извлечении квадратного корня учтены оба знака и четверть угла.
\item В сокращениях не допущено деление на нуль.
\end{itemize}

\clearpage
\section*{Ответы к тренировочному блоку}
\ruleline
\renewcommand{\arraystretch}{1.65}
\begin{center}
\begin{tabularx}{\linewidth}{@{}>{\centering\arraybackslash}p{12mm}X@{}}
\toprule
\textbf{№} & \textbf{Ответ}\\
\midrule
1 & $-\dfrac12$\\
2 & $135^\circ=\dfrac{3\pi}{4}$; $\cos135^\circ=-\dfrac{\sqrt2}{2}$\\
3 & $\dfrac12$\\
4 & $-\dfrac{\sqrt2}{2}$\\
5 & $2\sin7\alpha\cos7\alpha$\\
6 & $1-2\sin^24\alpha$\\
7 & $2{,}4$\\
8 & $14$\\
9 & $\dfrac35$\\
10 & $5$\\
\bottomrule
\end{tabularx}
\end{center}
\vspace{12mm}
\topic{Короткая памятка}
\begin{tabularx}{\linewidth}{@{}p{50mm}X@{}}
\toprule
\textbf{Что видите} & \textbf{Что применить}\\
\midrule
Сумму или разность углов & Формулу сложения для синуса или косинуса.\\
Аргументы $2x$ и $x$ & Формулу двойного угла.\\
Известен синус, требуется косинус & Основное тождество и знак по четверти.\\
Требуется тангенс & $\operatorname{tg}x=\sin x/\cos x$ при $\cos x\ne0$.\\
Большой угол с $\pi$ & Разложение на целое число $\pi$ или $2\pi$ и остаток.\\
Угол в градусах & Умножить число градусов на $\pi/180$.\\
\bottomrule
\end{tabularx}
\vfill
\begin{center}\small\color{muted} Лёвин Артём Александрович эксклюзивно для Марка Гукина\end{center}
\end{document}
